\section{Background and Motivation}

\subsection{Efficient Prefill}

Long-horizon and multi-turn agents make inference increasingly input-dominated.
Tool responses can add substantially more tokens than the actions that produced them.
Reducing attention cost alone does not eliminate the computation needed to propagate these tokens through the backbone.
Cross-layer KV-cache sharing can reduce both KV storage and prefill computation.
YOCO uses a self-decoder to construct a global KV cache shared across cross-decoder layers, allowing prefill cache construction to exit after the self-decoder~\citep{yoco}.
Gemma 3n adopts related cross-layer KV sharing to improve prefill efficiency~\citep{gemma3n}.
HySparse applies cross-layer sharing within its hybrid sparse-attention blocks: each full-attention layer shares its global KV cache and top-$k$ block indices with the following sparse-attention layers, though prefill still executes all layers~\citep{hysparse}.
HySparse2 further incorporates a YOCO-style structure into HySparse, combining reduced attention computation and KV storage with a shorter prefill path.

\subsection{KV Cache Compression}

Agentic workloads require more compact KV caches.
KV-cache compression spans four dimensions: head, sequence, layer, and precision.
Head-level methods share KV heads through GQA/MQA~\citep{gqa,mqa} or compress KV representations into latent states, as in MLA~\citep{deepseekv2}.
Sequence-level methods compress multiple tokens into fewer cache entries~\citep{deepseekv4}.
Layer-level methods share KV caches across layers~\citep{cla,yoco}.
Precision-level methods reduce the numerical precision of cached keys and values~\citep{liu2024kivi,hooper2024kvquant}.
Prior work focuses on intra-layer compression along the head, sequence, and precision axes.
HySparse2 instead pushes cross-layer sharing through KV Bridging and KV Reuse.

\subsection{Sparse Attention Granularity}

Sparse attention reduces long-context attention costs by restricting each query to a subset of KV entries~\citep{sparsetransformer}.
Its selection granularity creates an inherent trade-off between modeling accuracy and inference efficiency.
HySparse adopted block-level sparsity as a practical compromise, since evaluations showed no substantial accuracy disadvantage and the regular block structure enabled efficient kernels.
Modern agentic workloads, however, involve long-horizon multi-turn trajectories that require precise retrieval from long historical contexts.
Under these patterns, block-level selection exhibits a pronounced accuracy disadvantage, while token-level selection retrieves relevant evidence more faithfully. 
Recent advances in sparse kernels also make token-level implementations practical~\citep{tilelang}.
HySparse2 therefore adopts token-level sparsity for more precise context selection.
